\begin{figure}[H]
    \centering
    \begin{tikzpicture}[
    scale=1.5,
    cota/.style={<->, >=stealth, thin, black},
    chamada/.style={gray!80, thin},
    setaCarga/.style={->, >=stealth, black, thick},
    texto/.style={font=\small\sffamily}
]

    % --- Parâmetros de Geometria ---
    \def\vao{6.0}         % Comprimento do vão da laje
    \def\espPiso{0.25}    % Espessura da camada de piso
    \def\espLaje{0.80}    % Espessura da laje de concreto
    \def\espForro{0.15}   % Espessura da camada de forro

    % --- Coordenadas das Camadas (Origem na face inferior da laje) ---
    \coordinate (F_E_Inf) at (0, -\espForro);
    \coordinate (F_D_Sup) at (\vao, 0);
    
    \coordinate (L_E_Inf) at (0, 0);
    \coordinate (L_D_Sup) at (\vao, \espLaje);
    
    \coordinate (P_E_Inf) at (0, \espLaje);
    \coordinate (P_D_Sup) at (\vao, \espLaje + \espPiso);

    % --- 1. Camada de Forro / Acabamento Inferior ---
    \fill[gray!10!white] (F_E_Inf) rectangle (F_D_Sup);
    \draw[thick, black] (F_E_Inf) rectangle (F_D_Sup);

    % --- 2. Laje de Concreto Armado ---
    \fill[white] (L_E_Inf) rectangle (L_D_Sup);
    % Textura em hachura de concreto
    %\pattern[pattern=dots, pattern color=black!60] (L_E_Inf) rectangle (L_D_Sup);
    \draw[thick, black] (L_E_Inf) rectangle (L_D_Sup);

    % Representação das barras de armadura de aço (superior e inferior)
    %\draw[red!60!black, ultra thick] (0.2, 0.15) -- (\vao-0.2, 0.15);
    %\draw[red!60!black, ultra thick] (0.2, \espLaje-0.15) -- (\vao-0.2, \espLaje-0.15);

    % --- 3. Camada de Piso / Revestimento ---
    \fill[white] (P_E_Inf) rectangle (P_D_Sup);
    % Textura em hachura de concreto
    \pattern[pattern=vertical lines, pattern color=black!60]  (P_E_Inf) rectangle (P_D_Sup);
    \draw[thick, black] (P_E_Inf) rectangle (P_D_Sup);

    % --- Apóios da Laje nas extremidades (Vigas) ---
    %\draw[thick, fill=gray!50] (-0.5, -\espForro) rectangle (0, \espLaje);
    %\draw[thick, fill=gray!50] (\vao, -\espForro) rectangle (\vao+0.5, \espLaje);

    % --- 4. Carga de Uso Distribuída (Setas no Topo do Piso) ---
    \def\yPiso{\espLaje + \espPiso}
    \def\alturaSetas{0.7}
    \def\yTopoSetas{\yPiso + \alturaSetas}

    \draw[black, thick] (0, \yTopoSetas) -- (\vao, \yTopoSetas);

    \foreach \x in {0, 1, 2, 3, 4, 5, 6} {
        \draw[setaCarga] (\x, \yTopoSetas) -- (\x, \yPiso);
    }
    \node[texto, black, font=\bfseries\small, above] at (\vao/2, \yTopoSetas) {Carga de uso};

    % --- 5. Linhas de Chamada e Rótulos das Camadas ---
    \draw[chamada, ->, >=stealth] (\vao + 0.8, \yPiso - 0.12) -- (\vao - 0.5, \yPiso - 0.12);
    \node[texto, right] at (\vao + 0.8, \yPiso - 0.12) {Piso + Revestimento};

    \draw[chamada, ->, >=stealth] (\vao + 0.8, \espLaje/2) -- (\vao - 0.5, \espLaje/2);
    \node[texto, right] at (\vao + 0.8, \espLaje/2) {\begin{tabular}{l} Laje ($h=10\ cm$)) \end{tabular}};

    \draw[chamada, ->, >=stealth] (\vao + 0.8, -\espForro/2) -- (\vao - 0.5, -\espForro/2);
    \node[texto, right] at (\vao + 0.8, -\espForro/2) {Forro};

\end{tikzpicture}
    %\caption{Corte longitudinal da laje.}
    %\label{fig:corte-laje}
\end{figure}